% 逐元素差分在三元子集格 {1,2,3} 上的状态变化（累计值取第18.3节四元素设定在 {1,2,3} 的子集上的限制）。
% 初始存储 g：∅0, {1}2, {2}2, {3}2, {1,2}4, {1,3}7, {2,3}4, {1,2,3}9。
% 处理元素1后：{1,2}2, {1,3}5, {1,2,3}5（{1}仍为2，其余不变）。
% 处理元素2后：{1,2}0, {2,3}2, {1,2,3}0。
% 处理元素3后（即 h）：{1,3}3, {2,3}0，其余不变：∅0, 单项各2, {1,2}0, {1,2,3}0。
% 每一轮中，含被处理元素 i 的子集 S（黄色）减去 S\{i} 处的当前值；蓝色粗线连接 S\{i} 与 S，是这一轮读取的方向。
\begin{tikzpicture}[x=1mm,y=1mm,font=\figurefont\fontsize{8}{10}\selectfont,
  s/.style={draw=FigureStroke,fill=white,rounded corners=1.8pt,align=center,inner sep=1pt,minimum width=10mm,minimum height=7.2mm,line width=.7pt},
  e/.style={draw=FigureGrid,line width=.6pt},
  hot/.style={draw=FigureBlue,line width=1.3pt},
  ttl/.style={font=\figurefont\fontsize{8}{10}\selectfont,text=FigureInk}]
% ---------- 面板定义宏 ----------
% 参数：原点、本轮读取方向的连线、八个节点、标题
\newcommand{\drawpanel}[4]{%
  \begin{scope}[shift={#1}]
    \coordinate (pe) at (0,0);
    \coordinate (p1) at (-12,12); \coordinate (p2) at (0,12); \coordinate (p3) at (12,12);
    \coordinate (p12) at (-12,24); \coordinate (p13) at (0,24); \coordinate (p23) at (12,24);
    \coordinate (p123) at (0,36);
    \foreach \a/\b in {pe/p1,pe/p2,pe/p3,p1/p12,p1/p13,p2/p12,p2/p23,p3/p13,p3/p23,p12/p123,p13/p123,p23/p123}
      \draw[e] (\a)--(\b);
    #2
    #3
    \node[ttl] at (0,-8) {#4};
  \end{scope}
}
% 面板 (a)：初始
\drawpanel{(0,0)}{}{%
  \node[s] at (pe) {$\varnothing$\\$0$};
  \node[s] at (p1) {$\{1\}$\\$2$}; \node[s] at (p2) {$\{2\}$\\$2$}; \node[s] at (p3) {$\{3\}$\\$2$};
  \node[s] at (p12) {$\{1,2\}$\\$4$}; \node[s] at (p13) {$\{1,3\}$\\$7$}; \node[s] at (p23) {$\{2,3\}$\\$4$};
  \node[s] at (p123) {$\{1,2,3\}$\\$9$};
}{(a) 初始：存储$g$}
% 面板 (b)：处理元素1
\drawpanel{(38,0)}{%
  \foreach \a/\b in {pe/p1,p2/p12,p3/p13,p23/p123} \draw[hot] (\a)--(\b);
}{%
  \node[s] at (pe) {$\varnothing$\\$0$};
  \node[s,fill=FigureYellowFill] at (p1) {$\{1\}$\\$2$}; \node[s] at (p2) {$\{2\}$\\$2$}; \node[s] at (p3) {$\{3\}$\\$2$};
  \node[s,fill=FigureYellowFill] at (p12) {$\{1,2\}$\\$2$}; \node[s,fill=FigureYellowFill] at (p13) {$\{1,3\}$\\$5$}; \node[s] at (p23) {$\{2,3\}$\\$4$};
  \node[s,fill=FigureYellowFill] at (p123) {$\{1,2,3\}$\\$5$};
}{(b) 处理元素$1$后}
% 面板 (c)：处理元素2
\drawpanel{(76,0)}{%
  \foreach \a/\b in {pe/p2,p1/p12,p3/p23,p13/p123} \draw[hot] (\a)--(\b);
}{%
  \node[s] at (pe) {$\varnothing$\\$0$};
  \node[s] at (p1) {$\{1\}$\\$2$}; \node[s,fill=FigureYellowFill] at (p2) {$\{2\}$\\$2$}; \node[s] at (p3) {$\{3\}$\\$2$};
  \node[s,fill=FigureYellowFill] at (p12) {$\{1,2\}$\\$0$}; \node[s] at (p13) {$\{1,3\}$\\$5$}; \node[s,fill=FigureYellowFill] at (p23) {$\{2,3\}$\\$2$};
  \node[s,fill=FigureYellowFill] at (p123) {$\{1,2,3\}$\\$0$};
}{(c) 处理元素$2$后}
% 面板 (d)：处理元素3，得到 h
\drawpanel{(114,0)}{%
  \foreach \a/\b in {pe/p3,p1/p13,p2/p23,p12/p123} \draw[hot] (\a)--(\b);
}{%
  \node[s] at (pe) {$\varnothing$\\$0$};
  \node[s] at (p1) {$\{1\}$\\$2$}; \node[s] at (p2) {$\{2\}$\\$2$}; \node[s,fill=FigureYellowFill] at (p3) {$\{3\}$\\$2$};
  \node[s] at (p12) {$\{1,2\}$\\$0$}; \node[s,fill=FigureYellowFill] at (p13) {$\{1,3\}$\\$3$}; \node[s,fill=FigureYellowFill] at (p23) {$\{2,3\}$\\$0$};
  \node[s,fill=FigureYellowFill] at (p123) {$\{1,2,3\}$\\$0$};
}{(d) 处理元素$3$后，即$h$}
\end{tikzpicture}
